
\documentclass{../concepts/prepareprotokol} %všechny balíčky pro protokol + titulní strana
\usepackage{graphicx}
\usepackage{newfloat}
\usepackage{caption}
\usepackage{booktabs}
\usepackage{amsmath}
\usepackage{newfloat}
\usepackage{caption}
\usepackage{multirow}
\usepackage{siunitx}
\sisetup{inter-unit-product =$\cdot$}
\sisetup{separate-uncertainty=true}

\praktikum{II} %číslo praktika (I, II nebo III)
\autor{David Němec} %jméno autora
\datum{8.\,12.\,2025} %datum měření
\cislo{7} %číslo úlohy
\nazev{Měření indukčnosti a kapacity metodou přímou} %název úlohy

\DeclareFloatingEnvironment[fileext=los, placement={htbp}, name=Obrázek]{schema}

\newcommand{\nobars}{
Chybové úsečky nebyly pro svou malou velikost pro přehlednost vykresleny.}

\newcommand{\relative}[1]{
\frac{\sigma_{#1}}{#1}
}
\newcommand{\squared}[1]{
\left( #1 \right)^2
}
\newcommand{\relsqr}[1]{
\squared{\relative{#1}}
%\left( \frac{\sigma_{#1}}{#1} \right)^2
}

\begin{document}
\renewcommand{\figurename}{Graf}
\renewcommand{\refname}{Literatura}

\maketitle

% ========================== PRACOVNÍ ÚKOL ============================
\section{Pracovní úkol}

\begin{enumerate}
    \item Změřte velikost kapacit kondenzátorů z kapacitní dekády.
    \item Změřte závislost indukčnosti cívky na procházejícím proudu pro tyto případy:
    \begin{enumerate}
        \item cívka bez jádra (0-250 mA)
        \item cívka s otevřeným jádrem (0-250 mA)
        \item cívka s uzavřeným jádrem (0-250 mA)
    \end{enumerate}

    \item Změřte odpor cívky multimetrem a určete její kvalitu.
    \item Závislosti indukčnosti cívky na proudu měřené v úkolu 2 zakreslete do grafu přímo v praktiku.
\end{enumerate}

% ============================= TEORIE ================================
\section{Teorie}
\subsection{Odpor a kapacita kondenzátoru}

Reálný kondenzátor, který má kromě své kapacity $C$ i nenulový vlastní odpor $R_C$, lze popsat náhradním schématem (obr. \ref{sch:nahrkond}), kdy jsou paralelně zapojeny ideální kondenzátor $C$ a rezistor $R_C$. Kondenzátorem teče posuvný proud $i_p$ a rezistorem vodivostní proud $i_v$.

Pokud zapojíme reálný kondenzátor ke zdroji střídavého napětí o úhlové frekvenci $\omega$ a pro jeho odpor bude platit $R_C \gg (\omega C)^{-1}$, lze kapacitu kondenzátoru určit ze vztahu \cite{studijni}
\begin{equation} \label{eq:c}
    C = \frac{1}{2 \pi f} \frac{I}{U} \, ,
\end{equation}
kde $f$ je frekvence střídavého napětí, $I$ a $U$ jsou efektivní hodnoty proudu resp. napětí.

Při připojení ke stejnosměrnému napětí se kapacita kondenzátoru neprojeví a pro jeho odpor bude platit Ohmův zákon \cite{studijni}
\begin{equation} \label{eq:ohm}
    R_C = \frac{U}{I} \, .
\end{equation}

\begin{schema}[!htbp]
    \centering
    \includegraphics[width=0.5\linewidth]{uloha7/nahrkond.png}
    \caption{Náhradní schéma reálného kondenzátoru. \cite{studijni}}
    \label{sch:nahrkond}
\end{schema}

\begin{schema}[!htbp]
    \centering
    \includegraphics[width=0.5\linewidth]{uloha7/ckond.png}
    \caption{Zapojení obvodu při měření kapacity. \cite{pokyny}}
    \label{sch:ckond}
\end{schema}

\begin{schema}[!htbp]
    \centering
    \includegraphics[width=0.5\linewidth]{uloha7/rkond.png}
    \caption{Zapojení obvodu při měření odporu kondenzátoru. \cite{pokyny}}
    \label{sch:rkond}
\end{schema}

\subsection{Odpor a indukčnost cívky}
Reálnou cívku s indukčností $L$ a odporem $R_L$ lze popsat náhradním schématem (obr. \ref{sch:nahrcivka}), kdy jsou sériově zapojeny ideální cívka $L$ a rezistor $R_L$.

Pokud zapojíme cívku ke zdroji střídavého napětí o frekvenci $f$, lze indukčnost cívky vypočíst ze vztahu \cite{studijni}
\begin{equation} \label{eq:l}
    L = \frac{\sqrt{\frac{U^2}{I^2}-R_L^2}}{2 \pi f} \, .
\end{equation}

Měřítkem kvality cívky je její činitel jakosti, definovaný jako (upraveno podle \cite{studijni})
\begin{equation} \label{eq:q}
    Q = \frac{2 \pi f L}{R_L} \, .
\end{equation}

\begin{schema}[!htbp]
    \centering
    \includegraphics[width=0.5\linewidth]{uloha7/nahrcivka.png}
    \caption{Náhradní schéma reálné cívky. \cite{studijni}}
    \label{sch:nahrcivka}
\end{schema}

\begin{schema}[!htbp]
    \centering
    \includegraphics[width=0.5\linewidth]{uloha7/lcivka.png}
    \caption{Zapojení obvodu při měření cívky bez jádra a s otevřeným jádrem. \cite{pokyny}}
    \label{sch:lcivka}
\end{schema}

\begin{schema}[!htbp]
    \centering
    \includegraphics[width=0.5\linewidth]{uloha7/lcivka2.png}
    \caption{Zapojení obvodu při měření cívky s uzavřeným jádrem. \cite{pokyny}}
    \label{sch:lcivka2}
\end{schema}

\subsection{Čtyřbodové zapojení rezistoru} \label{four}

Pro přesnější měření odporu rezistoru (ale i cívky) lze využít čtyřbodového zapojení podle schématu na obrázku \ref{sch:four}. Pokud je voltmetr připojen k napěťovým svorkám $b$ a $b'$, nebude měření ovlivněno spádem napětí na přívodních vodičích mezi rezistorem a proudovými svorkami $a$ a $a'$. Za předpokladu nízkých hodnot $R$ oproti odporu voltmetru poteče přívodními vodiči k napěťovým svorkám zanedbatelný proudu a ani odpor těchto vodičů se neuplatní \cite{studijni}.

\begin{schema}
    \centering
    \includegraphics[width=0.5\linewidth]{uloha7/four.png}
    \caption{Čtyřbodové zapojení rezistoru. \cite{studijni}}
    \label{sch:four}
\end{schema}

\subsection{Metoda měření}

Nejprve byl obvod zapojen podle schématu \ref{sch:rkond} pro určení odporu kondenzátoru a odhad platnosti přepokladu vztahu (\ref{eq:c}) tj. $R_C \gg (\omega C)^{-1}$. Obvod byl připojen ke stejnosměrnému napětí, kdy se kapacita kondenzátoru neprojeví v jeho impedanci (frekvence je teoreticky nulová). 

Následně byl kondenzátor připojen ke střídavému napětí (odebíranému ze sítě s frekvencí \qty{50}{Hz}), které bylo transformátorem pro vyšší bezpečnost přeměněno na nižší hodnoty. Přístroj umožňující zapojování kondenzátorů do obvodu totiž nebyl chráněný vůči kontaktu.

Během všech měření s kondenzátory i cívkou bylo napětí nastaveno nejprve na nejvyšší hodnotu a následně bylo už jen snižováno, aby s jistotou nedošlo k překročení povolených hodnot proudů.

Při měření indukčnosti cívky bez jádra nebo s otevřeným jádrem byl opět jako zdroj použit transformátor snižující napětí, tentokrát ovšem z důvodu snažšího dosažení malých proudů. Při měření s cívkou s uzavřeným jádrem byl použit zdroj bez transformátoru, aby mohlo být pozorováno nasycení cívky při vyšších hodnotách proudu. Během měření s těmito vysokými hodnotami napětí bylo dbáno zvýšené opatrnosti.

Odpor cívky byl měřen čtyřbodovou metodou (schéma \ref{sch:four}), jelikož byla předpokládána jeho nízká hodnota. Za použití této medoty se totiž odpor přívodních vodičů neprojeví a výsledná hodnota nebude zatížena touto systematickou chybou.

\subsection{Měřicí přístroje a jejich nejistoty} \label{pristroje}
\begin{enumerate}
    \item Multimetr Fluke 175:
    
    Tímto multimetrem byly měřeny efektivní hodnoty střídavého napětí a proudu během měření kapacit kondenzátorů a také během měření indukčnosti cívky. Nejistota určení napětí je 1\,\% z měřené hodnoty + 3 digity a nejistota určení proudu je 1.5\,\% z měřené hodnoty + 3 digity.

    \item Multimetr Fluke 8808A:

    Tímto multimetrem byla ověřena frekvence střídavého napětí a dále také čtyřbodovým zapojením měřen odpor cívky. Nejistota určení frekvence je 0.01\,\% z měřené hodnoty + 0.003\,\% z rozsahu, nejistota určení odporu je 0.03\,\% z měřené hodnoty + 0.004\,\% z rozsahu.

    \item Analogový mikroampérmetr:

    Mikroampérmetrem byl měren proud při určování odporu kondenzátoru. Nejistota byla odhadnuta jako velikost nejmenšího dílku při daném rozsahu, tedy \qty{0.1}{\micro\ampere}.

\end{enumerate}


% ======================== VÝSLEDKY MĚŘENÍ ===========================
\section{Výsledky měření}

\subsection{Odpor a kapacita kondenzátorů}

Nejprve byla multimetrem Fluke 8808A ověřena hodnota frekvence střídavého napětí odebíraného ze sítě. Hodnota byla určena jako $f = \qty{50.000 \pm 0.005}{Hz}$, přičemž nejistota byla vzata jako nejistota přístroje.

Následně byl obvod zapojen podle schématu \ref{sch:rkond} a byl určen odpor kondenzátoru F s nejvyšší předpokládanou kapacitou, a tím i nejpřísnějším limitem na hodnotu odporu $R_C \gg (\omega C)^{-1}$. Efektivní hodnoty napětí a proudu byly změřeny jako $U= \qty{20.8\pm0.2}{V}$ resp. $I=\qty{5.0\pm0.1}{\micro\ampere}$. Jejich nejistoty byly vzaty jako nejistota přístroje. Odpor byl určen z Ohmova zákona (\ref{eq:ohm}) jako $R = \qty{4.16\pm0.10}{\mega\ohm}$, přičemž nejistota byla vypočtena podle \cite{englich} jako
\begin{equation}
    \sigma_{R_C} = R_C \sqrt{\relsqr{U} + \relsqr{I}} \, .
\end{equation}

Dále byl obvod zapojen podle schématu \ref{sch:ckond} a byla měřena závislost proudu na napětí pro připojené jednotlivé kondenzátory A, C, F a poté pro kombinaci C a F. Naměřené hodnoty ukazuje tabulka \ref{tab:kond}. Nejistoty napětí a proudu byly určeny jako nejistota měřidla (multimetru Fluke 175). Závislosti proudu na napětí pro jednotlivé kondenzátory jsou vykresleny v grafu \ref{fig:kond}. 

\begin{table}[!htbp]
\centering
\caption{Naměřené efektivní hodnoty napětí a proudu pro jednotlivé kondenzátory}
\begin{tabular}{cc|cc|cc|cc}
 \toprule
 \multicolumn{2}{c|}{kondenzátor A} & \multicolumn{2}{c|}{kondenzátor C} & \multicolumn{2}{c|}{kondenzátor F} & \multicolumn{2}{c}{kondenzátory C a F} \\
 \midrule
 $U_A\, (\mathrm{V})$ & $I_A\, (\mathrm{mA})$ & $U_C\, (\mathrm{V})$ & $I_C\, (\mathrm{mA})$ & $U_F\, (\mathrm{V})$ & $I_F\, (\mathrm{mA})$ & $U_{CF}\, (\mathrm{V})$ & $I_{CF}\, (\mathrm{mA})$ \\
 \midrule
$10.34 \pm 0.13$ & $2.15 \pm 0.06$ & $10.33 \pm 0.13$ & $4.50 \pm 0.10$ & $10.12 \pm 0.13$ & $10.28 \pm 0.18$ & $10.12 \pm 0.13$ & $14.4 \pm 0.2$ \\
 $8.51 \pm 0.12$ & $1.75 \pm 0.06$ & $8.45 \pm 0.11$ & $3.62 \pm 0.08$ & $8.60 \pm 0.12$ & $8.59 \pm 0.16$ & $8.66 \pm 0.12$ & $12.2 \pm 0.2$ \\
 $7.09 \pm 0.10$ & $1.43 \pm 0.05$ & $7.00 \pm 0.10$ & $2.96 \pm 0.07$ & $7.02 \pm 0.10$ & $6.98 \pm 0.13$ & $7.00 \pm 0.10$ & $9.73 \pm 0.18$ \\  
 $5.64 \pm 0.09$ & $1.14 \pm 0.05$ & $5.41 \pm 0.08$ & $2.27 \pm 0.06$ & $5.65 \pm 0.09$ & $5.54 \pm 0.11$ & $5.56 \pm 0.09$ & $7.75 \pm 0.15$ \\
 $3.62 \pm 0.04$ & $0.73 \pm 0.04$ & $3.52 \pm 0.04$ & $1.47 \pm 0.05$ & $3.55 \pm 0.04$ & $3.46 \pm 0.08$ & $3.58 \pm 0.04$ & $4.92 \pm 0.10$ \\
 $1.70 \pm 0.02$ & $0.34 \pm 0.04$ & $1.58 \pm 0.02$ & $0.65 \pm 0.04$ & $1.60 \pm 0.02$ & $1.52 \pm 0.05$ & $1.68 \pm 0.02$ & $2.28 \pm 0.06$ \\
 \bottomrule
\end{tabular}
\label{tab:kond}
\end{table}

\begin{figure}[hbtp]
    \centering
    \includegraphics[width=0.5\linewidth]{uloha7/kond.png}
    \caption{Závislost efektivní hodnoty proudu na napětí pro jednotlivé kondenzátory. Svislé chybové úsečky nebyly (kromě kombinace kondenzátorů C a F) pro svou malou velikost pro přehlednost vykresleny.}
    \label{fig:kond}
\end{figure}

Datové body proloženy přímkami $I = k U$ a kapacity jednotlivých kondenzátorů byly určeny lineární regresí ze směrnic $k$ těchto přímek. Podle rovnice (\ref{eq:c}) je totiž platí
\begin{equation} \label{eq:kc}
    k= 2 \pi f C \, .
\end{equation}

Takto určené kapacity kondenzátorů pak shrnuje tabulka \ref{tab:kapacity}. Nejistoty kapacit byly vypočtené podle \cite{englich} s využitím vztahu (\ref{eq:kc}) jako
\begin{equation}
    \sigma_C = C \relative{k} \, ,
\end{equation}
protože nejistota frekvence je oproti nejistotě směrnice zanedbatelná.

\begin{table}[!htbp]
\centering
\caption{Směrnice lineárních závislostí proudu na napětí a vypočtené kapacity kondenzátorů}
\begin{tabular}{c|cc}
 \toprule
 kondenzátor & $k\, (\mathrm{mA/V})$ & $C\, (\unit{\micro\farad})$ \\
 \midrule
 A & $0.204 \pm 0.003$ & $0.651 \pm 0.010$ \\
 C & $0.426 \pm 0.005$ & $1.355 \pm 0.015$ \\
 F & $0.992 \pm 0.008$ & $3.16 \pm 0.03$ \\
 C + F & $1.397 \pm 0.011$ & $4.45 \pm 0.04$ \\
 \bottomrule
\end{tabular}
\label{tab:kapacity}
\end{table}

\subsection{Odpor a indukčnost cívky}

Nejprve byl multimetrem Fluke 8808A čtyřbodovou metodou (viz kap. \ref{four} a schéma \ref{sch:four}) změřen odpor cívky bez jádra. Jeho hodnota je $R_L = \qty{2.721\pm0.009}{\ohm}$, přičemž nejistota byla vzata jako nejistota měřícího přístroje.

Následně byly pro tři případy (cívka bez jádra, s otevřeným jádrem, s uzavřeným jádrem) proměřeny závislosti indukčnosti cívky na procházejícím proudu. Obvod byl zapojen podle schématu \ref{sch:lcivka} pro cívku bez jádra a s otevřeným jádrem a podle schématu \ref{sch:lcivka2} pro cívku s uzavřeným jádrem. 

Multimetry byly měřeny efektivní hodnoty napětí a proudu, indukčnost z nich (a z odporu $R_L$) byla vypočtena podle vztahu (\ref{eq:l}). Naměřené hodnoty proudu a napětí a vypočtené hodnoty indukčnosti pro jednotlivé případy jsou v tabulkách \ref{tab:civka1}, \ref{tab:civka2} a \ref{tab:civka3}. Nejistoty proudu a napětí byly určeny ze specifikací přístroje (viz kap. \ref{pristroje}). Nejistota indukčnosti byla vypočtena podle \cite{englich} podle vztahu
\begin{equation}
    \sigma_L = \frac{1}{2 \pi f} \frac{1}{\sqrt{\frac{U^2}{I^2} - R_L^2}} \sqrt{\frac{U^2}{I^4}\sigma_U^2 + \frac{U^4}{I^6} \sigma_I^2 + {R_L}^2 \sigma_{R_L}^2} \, .
\end{equation}

\begin{table}[!hp]
\centering
\caption{Naměřené efektivní hodnoty proudu, napětí a vypočtené indukčnosti pro cívku bez jádra.}
\begin{tabular}{ccc}
 \toprule
 $I\, (\mathrm{mA})$ & $U\, (\mathrm{V})$ & $L\, (\mathrm{mH})$ \\
 \midrule
 $24.1 \pm 0.4$ & $0.1214 \pm 0.0015$ & $13.5 \pm 0.4$ \\
 $53.5 \pm 0.8$ & $0.270 \pm 0.003$ & $13.5 \pm 0.4$ \\
 $74.1 \pm 1.4$ & $0.373 \pm 0.004$ & $13.5 \pm 0.4$ \\
 $101.6 \pm 1.8$ & $0.511 \pm 0.005$ & $13.5 \pm 0.4$ \\
 $126 \pm 2$ & $0.637 \pm 0.009$ & $13.5 \pm 0.4$ \\
 $150 \pm 3$ & $0.754 \pm 0.011$ & $13.5 \pm 0.4$ \\
 $174 \pm 3$ & $0.875 \pm 0.012$ & $13.5 \pm 0.4$ \\
 $202 \pm 3$ & $1.018 \pm 0.013$ & $13.5 \pm 0.4$ \\
 $221 \pm 4$ & $1.111 \pm 0.014$ & $13.5 \pm 0.4$ \\
 $256 \pm 4$ & $1.288 \pm 0.016$ & $13.5 \pm 0.4$ \\
 \bottomrule
\end{tabular}
\label{tab:civka1}
\end{table}

\begin{table}[!hp]
\centering
\caption{Naměřené efektivní hodnoty proudu, napětí a vypočtené indukčnosti pro cívku s otevřeným jádrem.}
\begin{tabular}{ccc}
 \toprule
 $I\, (\mathrm{mA})$ & $U\, (\mathrm{V})$ & $L\, (\mathrm{mH})$ \\
 \midrule
 $22.1 \pm 0.4$ & $0.482 \pm 0.005$ & $69.0 \pm 1.4$ \\
 $42.3 \pm 0.7$ & $0.934 \pm 0.012$ & $69.7 \pm 1.5$ \\
 $74.2 \pm 1.4$ & $1.64 \pm 0.02$ & $69.6 \pm 1.6$ \\
 $100.0 \pm 1.8$ & $2.21 \pm 0.03$ & $69.8 \pm 1.5$ \\
 $126 \pm 2$ & $2.78 \pm 0.03$ & $69.9 \pm 1.5$ \\
 $149 \pm 3$ & $3.30 \pm 0.04$ & $70.0 \pm 1.4$ \\
 $175 \pm 3$ & $3.88 \pm 0.04$ & $70.1 \pm 1.4$ \\
 $200 \pm 3$ & $4.45 \pm 0.05$ & $70.3 \pm 1.4$ \\
 $228 \pm 4$ & $5.07 \pm 0.05$ & $70.3 \pm 1.4$ \\
 $251 \pm 4$ & $5.58 \pm 0.06$ & $70.4 \pm 1.4$ \\
 \bottomrule
\end{tabular}
\label{tab:civka2}
\end{table}

\begin{table}[!hp]
\centering
\caption{Naměřené efektivní hodnoty proudu, napětí a vypočtené indukčnosti pro cívku s uzavřeným jádrem.}
\begin{tabular}{ccc}
 \toprule
 $I\, (\mathrm{mA})$ & $U\, (\mathrm{V})$ & $L\, (\mathrm{H})$ \\
 \midrule
 $4.7 \pm 0.1$ & $1.13 \pm 0.01$ & $0.76 \pm 0.02$ \\
 $10.7 \pm 0.2$ & $3.55 \pm 0.04$ & $1.05 \pm 0.02$ \\
 $21.2 \pm 0.3$ & $10.56 \pm 0.14$ & $1.59 \pm 0.03$ \\
 $35.6 \pm 0.6$ & $23.7 \pm 0.3$ & $2.12 \pm 0.04$ \\
 $48.2 \pm 0.8$ & $35.1 \pm 0.4$ & $2.32 \pm 0.04$ \\
 $60.5 \pm 1.2$ & $45.1 \pm 0.5$ & $2.37 \pm 0.05$ \\
 $71.9 \pm 1.4$ & $53.5 \pm 0.6$ & $2.37 \pm 0.05$ \\
 $82.0 \pm 1.5$ & $60.6 \pm 0.9$ & $2.35 \pm 0.06$ \\
 $103.8 \pm 1.9$ & $73.6 \pm 1.0$ & $2.26 \pm 0.05$ \\
 $126 \pm 2$ & $85.2 \pm 1.2$ & $2.14 \pm 0.05$ \\
 $148 \pm 3$ & $94.6 \pm 1.2$ & $2.03 \pm 0.04$ \\
 $170 \pm 3$ & $103.2 \pm 1.3$ & $1.93 \pm 0.04$ \\
 $192 \pm 3$ & $110.9 \pm 1.4$ & $1.84 \pm 0.04$ \\
 $216 \pm 4$ & $118.1 \pm 1.5$ & $1.74 \pm 0.04$ \\
 \bottomrule
\end{tabular}
\label{tab:civka3}
\end{table}

Výsledné závislosti indukčnosti na proudu pro jednotlivé případy jsou vykresleny v grafech \ref{fig:civka1}, \ref{fig:civka2} a \ref{fig:civka3}.

\begin{figure}[!hp]
    \centering
    \includegraphics[width=0.5\linewidth]{uloha7/fig1.png}
    \caption{Závislost indukčnosti cívky na proudu pro cívku bez jádra. Vodorovné chybové úsečky nebyly pro svou malou velikost pro přehlednost vykresleny.}
    \label{fig:civka1}
\end{figure}

\begin{figure}[!hp]
    \centering
    \includegraphics[width=0.5\linewidth]{uloha7/fig2.png}
    \caption{Závislost indukčnosti cívky na proudu pro cívku s otevřeným jádrem.}
    \label{fig:civka2}
\end{figure}

\begin{figure}[!htbp]
    \centering
    \includegraphics[width=0.5\linewidth]{uloha7/fig3.png}
    \caption{Závislost indukčnosti cívky na proudu pro cívku s uzavřeným jádrem.}
    \label{fig:civka3}
\end{figure}

Činitel jakosti cívky byl vypočten podle vztahu (\ref{eq:q}), přičemž jeho nejistota byla určena podle \cite{englich} jako (nejistota frekvence je zanedbatelná oproti nejistotě indukčnosti nebo odporu)
\begin{equation}
    \sigma_Q = Q \sqrt{\relsqr{L} + \relsqr{R_L}} \, .
\end{equation}

Pro cívku bez jádra a s otevřeným jádrem byl činitel jakosti vypočten ze střední hodnoty indukčnosti, pro cívku s uzavřeným jádrem byla použita její maximální hodnota $L_{max} = \qty{2.37\pm0.05}{\henry}$ v bodě s proudem \\ $I_{max}= \qty{65\pm10}{\milli\ampere}$. Nejistota proudu v maximu indukčnosti byla pouze odhadnuta z grafu \ref{fig:civka3}.
Přehled indukčností a činitelů jakosti cívky pro jednotlivé případy je v tabulce \ref{tab:kvality}.

\begin{table}
    \centering
    \caption{Indukčnosti a činitele jakosti cívky bez jádra, s otevřeným jádrem a s uzavřeným jádrem. Pro první dva případy je uvedena střední hodnota indukčnosti, pro cívku s uzavřeným jádrem maximální hodnota indukčnosti.}
    \begin{tabular}{c|c|c}
        \toprule
         cívka & $L \, (\unit{\milli\henry})$ & $Q$ \\
         \midrule
         bez jádra & $13.5\pm0.4$ & $1.56\pm0.05$ \\
         s otevřeným jádrem & $69.9\pm1.4$ & $8.07\pm0.16$ \\
         s uzavřeným jádrem & $2370\pm50$ & $274\pm6$ \\
         \bottomrule
    \end{tabular}
    \label{tab:kvality}
\end{table}

% ============================ DISKUZE ===============================
\section{Diskuze}

Při měření odporu kondenzátoru byl obvod zapojen podle schématu \ref{sch:rkond}, tedy s voltmetrem před ampérmetrem. Ampérmetr díky tomu měří proud přímo na kondenzátoru, nikoliv celkový proud kondenzátorem i voltmetrem. Odpor ampérmetru je navíc oproti odporu kondenzátoru zanedbatelný, takže i voltmetr ukazuje správné hodnoty. Odpor byl změřen jako $R_C = \qty{4.16\pm0.10}{\mega\ohm}$, s relativní nejistotou 2.4\,\%. Ta je daná především nejistotou určení proudu, tedy nejistotou analogového mikroampérmetru.

Kapacity jednotlivých kondenzátorů A, C, F a kombinace C a F byly určeny lineární regresí ze závisloti proudu na napětí jako $C_A = \qty{0.651 \pm 0.010}{\micro\farad}$, $C_C = \qty{1.355 \pm 0.015}{\micro\farad}$, $C_F = \qty{3.16 \pm 0.03}{\micro\farad}$ a \\$C_{CF} = \qty{4.45 \pm 0.04}{\micro\farad}$, s relativními nejistotami okolo 1\,\%. Nejistota je daná nejistotou směrnice lineární závislosti proudu na napětí. 

Hodnotu kapacity kombinace kondenzátorů C a F lze ověřit výpočtem. Přístroj umožnoval zapojovat kondenzátory do obvodu paralelně, tedy celková kapacita bude součtem kapacit jednotlivých kondenzátorů. Ta tedy vychází $C'_{CF} = \qty{4.51\pm0.03}{\micro\farad}$, přičemž nejistota byla vypočtena jako $\sigma_{C'_{CF}} = \sqrt{\sigma_{C_C}^2 + \sigma_{C_F}^2}$. V rámci nejistoty jsou tedy obě hodnoty $C_{CF}$ a $C'_{CF}$ v souladu.

I pro nejvyšší kapacitu $C_{CF} = \qty{4.45 \pm 0.04}{\micro\farad}$ použitou během měření je splněn předpoklad \\$R_C \gg (2 \pi f C)^{-1}$, neboť hodnota $(2 \pi f C_{CF})^{-1}$ vychází $\qty{715\pm6}{\ohm}$, což je téměř 6000krát menší než \\$\qty{4.16\pm0.10}{\mega\ohm}$. Odpor byl měřen jen pro kondenzátor F, pro ostatní kondenzátory se tedy může jeho hodnota lišit. Předpoklad platnosti rovnice (\ref{eq:c}) však bude splněn, i pokud si budou odpory odpovídat pouze řádově. Vztah (\ref{eq:c}) pro výpočet kapacity lze tedy bezpečně použít pro všechny kondenzátory.

Při měření indukčnosti cívky byl obvod zapojen podle schématu \ref{sch:lcivka}, tedy s voltmetrem za ampérmetrem. Odpor cívky je totiž malý a úbytek napětí na ampérmetru by mohl systematicky ovlivnit měření. Oproti voltmetru je však nyní odpor cívky zanedbatelný, tedy celkový proud ampérmetrem se velmi blíží hodnotě proudu přímo na cívce. Odpor byl změřen jako $R_L = \qty{2.721\pm0.009}{\ohm}$ s relativní nejistotou jen 0.3\,\% díky velmi přesnému přístroji Fluke 8808A. Díky použité čtyřbodové metodě měření odporu není tato hodnota zatížena systematickou chybou danou odporem přívodních vodičů.

Z grafu \ref{fig:civka1} bylo ověřeno, že indukčnost cívky bez jádra nezávisí na procházejícím proudu. Hodnoty jsou v rámci nejistoty skutečně konstantní. Jejich střední hodnota vychází $L = \qty{13.5\pm0.4}{\milli\henry}$. Kvalita (činitel jakosti) cívky vypočtený z této střední hodnoty pak je $Q = 1.56\pm0.05$. Nejistoty indukčnosti a činitele jakosti jsou dané především nejistotou měření proudu.

Hodnoty indukčnosti cívky s otevřeným jádrem podle grafu \ref{fig:civka2} s rostoucím proudem velmi mírně rostou, nejistota je však příliš velká na to, aby bylo možné z toho vyvozovat nějaké závěry. Pro ověření této stoupající tendence indukčnosti by bylo třeba přesnější měření (především hodnot proudu, které mají kvůli své velké relativní nejistotě největší vliv na nejistotu indukčnosti). Střední hodnota indukčnosti cívky s otevřeným jádrem vychází $L = \qty{69.9\pm1.4}{\milli\henry}$, čemuž odpovídá činitel jakosti $Q = 8.07\pm0.16$.

Nakonec byla proměřována závislost indukčnosti na proudu pro cívku s uzavřeným jádrem. Bylo pozorováno (viz graf \ref{fig:civka3}), že indukčnost s rostoucím proudem rostla až do maximální hodnoty $L_{max} = \qty{2.37\pm0.05}{\henry}$, která odpovídala proudu přibližně $I_{max}= \qty{65\pm10}{\milli\ampere}$. V tomto bodě došlo k nasycení jádra a indukčnost s dalším nárůstem proudu pozvolna klesala.

Činitel jakosti je největší (a to dokonce řádově) pro cívku s uzavřeným jádrem, protože v tomto případě má cívka díky jádru největší hodnotu indukčnosti. Odpor cívky byl uvažován stejný pro všechny případy. Činitel jakosti je však v tomto případě uveden pouze pro maximální hodnotu indukčnosti, pro jiné hodnoty proudu bude jeho velikost menší. Přesto je pro všechny body z rozsahu měření nejkvalitnější cívka s uzavřeným jádrem.

% ============================= ZÁVĚR ================================
\section{Závěr}

\hspace{\parindent} Odpor kondenzátoru F byl změřen jako $R_C = \qty{4.16\pm0.10}{\mega\ohm}$.
Kapacita kondenzátoru A byla určena jako \\$C_A = \qty{0.651 \pm 0.010}{\micro\farad}$,
kondenzátoru C jako $C_C = \qty{1.355 \pm 0.015}{\micro\farad}$,
kondenzátoru F jako \\$ C_F = \qty{3.16 \pm 0.03}{\micro\farad}$,
a paralelní kombinace kondenzátorů C a F jako $C_{CF} = \qty{4.45 \pm 0.04}{\micro\farad}$.
\vspace{1mm}

Odpor cívky byl změřen jako $R_L = \qty{2.721\pm0.009}{\ohm}$.
Indukčnost cívky bez jádra byla určena jako \\$L = \qty{13.5\pm0.4}{\milli\henry}$, přičemž její činitel jakosti vychází $Q = 1.56\pm0.05$,
indukčnost cívky s~otevřeným jádrem jako $L = \qty{69.9\pm1.4}{\milli\henry}$ s činitelem jakosti $Q = 8.07\pm0.16$
a indukčnost cívky s~uzavřeným jádrem byla určena jako $L = \qty{2.37\pm0.05}{\henry}$ s činitelem jakosti $Q = 274\pm6$.


% =========================== LITERATURA =============================
\begin{thebibliography}{9}

\bibitem{studijni} Kolektiv ZFP KVOF MFF UK: Měření indukčnosti a kapacity metodou přímou - studijní text [online]. [cit. 14.12.2025], dostupné z https://physics.mff.cuni.cz/vyuka/zfp/\_media/zadani/texty/txt\_207.pdf

\bibitem{pokyny} Kolektiv ZFP KVOF MFF UK: Měření indukčnosti a kapacity metodou přímou - pokyny k měření [online]. [cit. 14.12.2025], dostupné \\
z https://physics.mff.cuni.cz/vyuka/zfp/\_media/zadani/pokyny/mereni\_207.pdf

\bibitem{englich} J. Englich: Úvod do praktické fyziky I: Zpracování výsledků měření. 1. vyd. Praha: Matfyzpress, 2006

\end{thebibliography}

\end{document}